
\documentclass{article} % For LaTeX2e
\usepackage{iclr2025_conference,times}

% Optional math commands from https://github.com/goodfeli/dlbook_notation.
\input{math_commands.tex}
\input{header.tex}

\usepackage{hyperref}
\hypersetup{
	final,
	colorlinks,
	linkcolor=ourred,
	citecolor=ourblue,
	urlcolor=url,
}
\usepackage{url}
\usepackage{booktabs}
% \usepackage{todonotes}
\usepackage{caption}
\usepackage{makecell}
\usepackage{adjustbox} % To adjust table size
\usepackage{wrapfig}
\usepackage{amsmath, amssymb}

\newcommand{\guanghan}[1]{\textcolor{purple}{[Guanghan: #1]}}
\newcommand{\hao}[1]{\textcolor{red}{[Hao: #1]}}
\newcommand{\vkl}[1]{\textcolor{blue}{[VK: #1]}}



\title{Simple Guidance Mechanisms\\for Discrete Diffusion Models}

% Authors must not appear in the submitted version. They should be hidden
% as long as the \iclrfinalcopy macro remains commented out below.
% Non-anonymous submissions will be rejected without review.

\author{Yair Schiff\thanks{Corresponding author. Work done while at InstaDeep.~$^\dagger$Equal contribution.}~~$^\dagger$, Subham Sekhar Sahoo$^\dagger$,  Hao Phung$^\dagger$, Guanghan Wang$^\dagger$,
\\
\textbf{Alexander Rush, \& Volodymyr Kuleshov}\\
Department of Computer Science, Cornell University\\
\texttt{\{yzs2,sss284,gw354,htp26,amr459,vk379\}@cornell.edu} \\
\AND
Hugo Dalla-Torre, Sam Boshar, Bernardo P. de Almeida, \& Thomas Pierrot\\
InstaDeep\\
\texttt{\{h.dalla-torre,s.boshar,b.dealmeida,t.pierrot\}@instadeep.com}
}
\newcommand{\fix}{\marginpar{FIX}}
\newcommand{\new}{\marginpar{NEW}}

\iclrfinalcopy % Uncomment for camera-ready version, but NOT for submission.
% \iclrpreprintcopy
\begin{document}


\maketitle

\begin{abstract}
Diffusion models for continuous data gained widespread adoption owing to their high quality generation and control mechanisms. However, controllable diffusion on discrete data faces challenges given that continuous guidance methods do not directly apply to discrete diffusion. Here, we provide a straightforward derivation of classifier-free and classifier-based guidance for discrete diffusion, as well as a new class of diffusion models that leverage uniform noise and that are more guidable because they can continuously edit their outputs. We improve the quality of these models with a novel continuous-time variational lower bound that yields state-of-the-art performance, especially in settings involving guidance or fast generation. Empirically, we demonstrate that our guidance mechanisms combined with uniform noise diffusion improve controllable generation relative to autoregressive and diffusion baselines on several discrete data domains, including genomic sequences, small molecule design, and discretized image generation. Code to reproduce our experiments is available \href{https://github.com/kuleshov-group/discrete-diffusion-guidance}{here}.
\end{abstract}

\section{Introduction}\label{sec:intro}

Diffusion models \citep{sohl2015deep,ho2020denoising} gained widespread adoption in image generation and signal processing in part due to their high controllability via mechanisms such as classifier-based \citep{dhariwal2021agn} and classifier-free guidance \citep{nichol2021glide, ho2022classifier}.
Tasks where guidance plays a key role include MRI denoising \citep{song2019generative}, 3D reconstruction \citep{poole2022dreamfusion,gao2024cat3dcreate3dmultiview}, and conditional generation \citep{saharia2022photorealistic, gokaslan2024commoncanvas}.

However, applying controllable diffusion-based generation to tasks where the data is discrete (e.g., molecule design or text generation) presents challenges.
First, standard diffusion models and their guidance mechanisms are not directly applicable, since they require taking gradients with respect to the data, and these are not defined in discrete settings. Second, popular discrete extensions of diffusion \citep{sahoo2024simple, shi2024simplified} cannot perform multiple editing passes on generated tokens, hence are not ideal for controllable generation. Third, the performance of discrete diffusion models (measured by perplexity) lags behind autoregressive (AR) models, especially for classes of diffusion that are amenable to control, such as uniform noise \citep{austin2021structured,lou2023discrete}. 

Here, we propose discrete diffusion models and guidance mechanisms that are effective at controllable generation and that address the above challenges. 
First, we provide straightforward and easy-to-implement adaptations of classifier-based and classifier-free guidance for discrete diffusion models.
Second, we revisit uniform noise diffusion language models (UDLM), which undo random token perturbations and are particularly amenable to guidance, since they can repeatedly edit their samples \citep{austin2021structured} and thus correct errors.
We address performance issues that plagued previous iterations of uniform noise discrete diffusion by deriving a continuous-time version of the evidence lower bound, which tightens the variational gap \citep{kingma2021variational,sahoo2024simple}.

We demonstrate the effectiveness of guidance with discrete diffusion models on several domains: genomics, molecule strings, and discretized images.
We find that discrete diffusion models are more controllable than AR when paired with classifier-free guidance.
Moreover, our proposed classifier-based method improves upon previous guidance mechanisms for diffusion \citep{gruver2024protein} and AR \citep{dathathri2019plug, yang2021fudge}.
Our language modeling experiments also reveal that contrary to a widely-held belief \citep{austin2021structured,lou2023discrete}, uniform noise diffusion can attain state-of-the-art performance on small vocabulary datasets (e.g,. molecules, DNA) and that UDLM attains a new state-of-the-art in perplexity among uniform noise diffusion models.

In summary, our contributions are as follows:
\begin{itemize}
    \item We provide simple and effective discrete classifier-based and classifier-free guidance.
    \item We introduce UDLM, a class of discrete diffusion models particularly amenable to guidance, for which we derive a tightened ELBO that significantly improves their performance.
    \item Across three domains, we demonstrate that discrete guidance yields better controllable generation compared to strong AR baselines and previous diffusion guidance methods.
\end{itemize}

\begin{figure}
    \centering
    \includegraphics[width=0.9\linewidth]{images/graphical_abstract.pdf}
    \caption{\textit{(Left)} Adapting guidance to discrete diffusion.
    Models output a factorized discrete distribution for each denoised token.
    With our guidance mechanisms, we adjust these probabilities according to a guidance model -- either a conditional diffusion model in classifier-free guidance (\Secref{subsec:cfg}) or a separately trained classifier for classifier-based guidance (\Secref{subsec:cbg}).
    \textit{(Right)} Relative to autoregressive models, which make local predictions one token at a time, discrete diffusion models denoise the entire sequence at every iteration, allowing for more guidable outputs.}
    \label{fig:graphical_abstract}
\vspace{-1em}
\end{figure}

\section{Background}\label{sec:background}

\textbf{Notation}~
Let $\onehots$ be the space of all one-hot tokens over some vocabulary consisting of $\vocabsize$ unique characters: $\onehots = \{\z \in \{0, 1\}^{\vocabsize} : \sum_i \z_i = 1\}\subset \Delta^\vocabsize$, where $\Delta^N$ represents the simplex over $N$ categories.
Let $\vone$ be a $\vocabsize$-dimensional column vector of all ones, and denote the Hadamard product between two vectors as $\odot.$
We define $\z^{(1:L)}$ as a sequence of $L$ tokens, where $\z^{(\ell)} \in \onehots,$ for all tokens $\ell \in \{1, \ldots, L\},$ and use $\onehots^L$ to denote the set of all such sequences.
Finally, let $\cat(\cdot; p)$ denote the categorical distribution with probability vector $p \in \Delta^\vocabsize.$

\subsection{Discrete Diffusion Models}\label{subsec:background_diffusion}

Diffusion models are a class of generative models defined by a denoising network $p_\theta$ that is trained to remove noise from latent variables $\z_t$.
These latents are generated by a fixed corruption process $q$ that, starting from clean data $\x$ drawn from the data distribution $q(\x)$, increasingly adds more noise to $\z_t$, as $t$ increases \citep{sohl2015deep, 
song2019generative, ho2020denoising}.

In discrete denoising diffusion probabilistic models (D3PM; \citet{austin2021structured}), the noising process is defined in terms
of a transition matrix $Q_{t|s}$ 
whose $(i,j)^{\text{th}}$ entry is the probability of transitioning to the $i$-th state at time $t$ given the $j$-th state at time $s$.
This induces a Markov corruption process where we have $q(\z_t \mid \z_s) = \cat(\z_t; Q_{t|s}\z_s)$.

\citet{sahoo2024simple} build off this framework to introduce specialized algorithms that are both simpler and more effective than the general D3PM framework.
They focus on a specific class of forward processes from D3PM that can be defined as interpolations between clean data and a noisy prior $\prior$, and we adopt their notation below:
\begin{align}\label{eq:interpolating_forward}
    q(\z_t | \x) = \cat(\z_t; \at \x + (1 - \at) \prior{}),
\end{align}
where $\at = \alpha(t)$ is a noise schedule monotonically decreasing in $t.$
Defining $\ats = \at / \as,$ this class of processes admit the following posteriors
\begin{equation}\label{eq:interpolating_posterior}
    q(\z_s | \z_t, \x) = \cat \left(
        \z_s; \frac{[\ats \z_t + (1 - \ats)\vone \prior^\top \z_t]  \odot [\as \x + (1 - \as) \prior]}{
             \at \z_t^\top \x  + (1 - \at)\z_t^\top\prior} \right).
\end{equation}

Of note, for absorbing-state diffusion, where $\prior = \mask,$ a one-hot vector at the special $\mask$ token index, \cite{sahoo2024simple} show that if $\z_t \neq \mask$ then $q(\z_s | \z_t, \x) = \cat (\z_s; \z_t)$, which reflects the fact that unasked tokens at time $t$ must remain unmasked for all time $s < t$.

\subsection{Diffusion Guidance}\label{subsec:background_guidance}
For continuous data, diffusion models have demonstrated state-of-the-art controllable generation by means of classifier-based \citep{sohl2015deep, dhariwal2021agn} and classifier-free guidance \citep{nichol2021glide, ho2022classifier, saharia2022photorealistic}. These approaches rely on different ways of expressing the score of a distribution conditioned on $y$.

\textbf{Classifier-based}~
Classifier-based generation employs a diffusion model to iteratively sample from a tempered distribution $p^{\gamma}(\z_s \mid y, \z_t) \propto p(y \mid \z_s)^{\gamma}p_\theta(\z_s \mid\z_t),$ where $\gamma$ represents an inverse temperature parameter, $p_\theta(\z_s \mid\z_t)$ is a pre-trained diffusion model, and $p(y \mid \z_s)$ is a classifier:
\begin{align}\label{eq:classifier_based}
    \nabla_{\z_s} \log p^\gamma(\z_s \mid y, \z_t) = \gamma\nabla_{\z_s}\log p(y \mid \z_s) + \nabla_{\z_s}\log p_\theta(\z_s \mid \z_t).
\end{align}

\textbf{Classifier-free}~
We can also observe that appyling Bayes' rule to $p(y \mid \x)$ and differentiating with respect to the input yields $\nabla_\x \log p(y \mid x) = \nabla_\x \log p(x \mid y) - \nabla_\x \log p(x).$
Applying this to $p(y \mid \z_s)$ and plugging into \eqref{eq:classifier_based} gives us the formulation for classifier-free guidance:
\begin{align}
    \nabla_{\z_s} \log p^{\gamma}(\z_s \mid y, \z_t) &= \gamma\cdot[\nabla_{\z_s}\log p_\theta(\z_s \mid y, \z_t) - \nabla_{\z_s} \log p_\theta(\z_s \mid \z_t)] + \nabla_{\z_s} \log p_\theta(\z_s \mid \z_t) \nonumber \\
    &= \gamma\nabla_{\z_s}\log p_\theta(\z_s \mid y, \z_t) + (1-\gamma) \nabla_{\z_s} \log p_\theta(\z_s \mid \z_t). \label{eq:classifier_free}
\end{align}
with $p_\theta(\z_s \mid y, \z_t)$ conditionally and $p_\theta(\z_s \mid \z_t)$ unconditionally trained diffusion models.
% In practice, this is often implemented by using a single model $p_\theta$ and randomly dropping out / masking the conditioner $y$ during training.

\section{Guidance Algorithms for Discrete Diffusion}\label{sec:guidance}
Applying guidance to discrete diffusion is challenging because guidance terms are not differentiable with respect to the discrete representations $\z_t$ \citep{gruver2024protein}.
Here, we introduce two guidance algorithms for discrete diffusion that circumvent the issue of non-differentiability.
As in \Secref{subsec:background_guidance}, we formalize the guidance term as a probability $p(y | \z) \in \Delta^K$, where $y \in \{1, \ldots, K\}$ is one of $K$ possible classes, and we begin with the distribution that controls the strength of the guidance term via the $\gamma$ temperature parameter (with $c$ denoting the logarithm of the normalization constant):
\begin{equation}\label{eq:guidance_unormalized}
    \log p^\gamma(\z_s \mid y, \z_t) = \gamma\log p(y\mid\z_s, \z_t) + \log p(\z_s \mid \z_t) + c.
\end{equation}

\subsection{Classifier-free Guidance}\label{subsec:cfg}
\textbf{Guidance for a Single Token}~
To derive classifier-free guidance, we apply Bayes' rule to the first term on the right-hand side of \eqref{eq:guidance_unormalized}:
\begin{equation*}
    \log p^\gamma(\z_s \mid y, \z_t) = \gamma\log p(\z_s \mid y, \z_t) - \gamma\log p(\z_s \mid \z_t) + \gamma\log p(y\mid \z_t) + \log p(\z_s \mid \z_t) + c. 
\end{equation*}
Absorbing $\gamma\log p(y\mid \z_t)$ into the constant $c$ and grouping terms yields:
\begin{equation}\label{eq:cfg_single_token}
    \log p^\gamma(\z_s \mid y, \z_t) = \gamma\log p(\z_s \mid y, \z_t) + (1-\gamma)\log p(\z_s \mid \z_t) + c.
\end{equation}

\textbf{Extension to Sequences}~
In the derivation for \eqref{eq:cfg_single_token}, we can replace the single-token latent variables with sequences of multiple tokens:
\begin{align}\label{eq:cfg_sequences}
\begin{split}
    \log p^\gamma(\z_s^{(1:L)} \mid y, \z_t^{(1:L)}) &= \gamma\log p(\z_s^{(1:L)} \mid y, \z_t^{(1:L)}) + (1-\gamma)\log p(\z_s^{(1:L)} \mid \z_t^{(1:L)}) + c,\\
    \implies p^\gamma(\z_s^{(1:L)} \mid y, \z_t^{(1:L)}) &\propto p(\z_s^{(1:L)} \mid y, \z_t^{(1:L)})^{\gamma}p(\z_s^{(1:L)} \mid \z_t^{(1:L)})^{(1-\gamma)}
\end{split}
\end{align}

\textbf{Discrete Classifier-Free Guidance}~
We can model the probabilities in \eqref{eq:cfg_sequences} by training a conditional denoising diffusion network $p_\theta(\z_s^{(1:L)} \mid y, \z_t^{(1:L)})$ and an unconditional one $p_\theta(\z_s^{(1:L)} \mid \z_t^{(1:L)})$.
In practice, we train these models in tandem by randomly `dropping out' or `masking' the conditioner during training to simulate the unconditional diffusion model.
Each of these distributions is thus modeled to factorize independently along the sequence length dimension of $\z_s^{(1:L)}$ when conditioned on $\z_t^{(1:L)}$, allowing us to efficiently sample from the tempered guidance distribution as follows:
\begin{align}
    p^\gamma_\theta(\z_s^{(1:L)} \mid \z_t^{(1:L)}, y) = \prod_{\ell=1}^{L}\frac{1}{Z^{(\ell)}}p_\theta(\z_s^{(\ell)} \mid y, \z_t^{(1:L)})^\gamma p_\theta(\z_s^{(\ell)}\mid\z_t^{(1:L)})^{(1-\gamma)},
\end{align}
where $Z^{(\ell)} = \sum_{\z_s'}p_\theta(\z_s' \mid \z_t^{(1:L)}, y)^\gamma p_\theta(\z_s'\mid\z_t^{(1:L)})^{(1-\gamma)}$ is the per-token partition function.

We refer to this method as \textbf{D-CFG} for \textbf{D}iscrete \textbf{C}lassifier-\textbf{F}ree \textbf{G}uidance.


\subsection{Classifier-based Guidance}\label{subsec:cbg}

\textbf{Guidance for a Single Token}~
Exponentiating \eqref{eq:guidance_unormalized} lets us draw classifier-guided samples from the following tempered distribution:
\begin{equation}\label{eq:cbg_single_token}
    p^\gamma(\z_s \mid \z_t, y) = \frac{p(y \mid \z_s, \z_t)^\gamma p(\z_s \mid \z_t)}{\sum_{\z_s'}p(y \mid \z_s', \z_t)^\gamma p(\z_s' \mid \z_t)},
\end{equation}
which we can tractably normalize, as we only sum over the $N$ possible values of a single token $\z_s.$

\textbf{Extension to Sequences}~
Unfortunately, extending \eqref{eq:cbg_single_token} to sequences is not trivial, since the classifier $p(y \mid \z_s^{(1:L)}, \z_t^{(1:L)})$ does not necessarily factorize independently across the sequence, which would potentially force us to compute a normalization constant with $N^{L}$ terms, accounting for every possible value of $\z_s^{(1:L)}.$
To alleviate this issue, we require the assumption that conditioned on $\z_t^{(1:L)}$, the tempered distribution $p^\gamma(\z_s^{(1:L)} \mid \z_t^{1:L}, y)$ factorizes independently across tokens.
Therefore, we can focus on the tempered distirbution of each token $\z_s^{(\ell)},$ for $\ell \in 1,\ldots, L$:
\begin{equation}\label{eq:cbg_sequences}
    p^\gamma(\z_s^{(\ell)} \mid \z_t^{(1:L)}, y) \propto p(y \mid \z_s^{(\ell)}, \z_t^{(1:L)})^\gamma p(\z_s^{(\ell)} \mid \z_t^{(1:L)}).
\end{equation}

\textbf{Discrete Classifier-Based Guidance}~
We first introduce the following additional notation: given $\z^{(1:L)}$, let $\Tilde{\mathcal{Z}}_\ell(\z^{(1:L)})$ be the set of sequences $\Tilde{\z}^{(1:L)}$ for which $\Tilde{\z}^{(\ell')} = \z^{(\ell')}$ for all $\ell' \neq \ell,$ i.e., the set of sequences that are either the same as or only differ in position $\ell$ relative to $\z^{(1:L)}.$
We can sample from $p^\gamma(\z_s^{(\ell)} \mid \z_t^{(1:L)}, y)$ by training a classifier $p_\phi: \onehots^L \rightarrow \Delta^K$ on noised latents $\z_t^{(1:L)}$ for $t \in [0, 1]$ and use this to model the first term on the right-hand side of \eqref{eq:cbg_sequences} by only evaluating $p_\phi$ on sequences for which $\z_s^{(1:L)}$ and $\z_t^{(1:L)}$ differ by at most the token at position $\ell$:
\begin{align*}
    p(y \mid \z_s^{(\ell)}, \z_t^{(1:L)}) \approx p_\phi(y \mid \Tilde{\z}^{(1:L)}), &&
    \text{for ~~}\Tilde{\z}^{(1:L)} = \Big[\z_t^{(1:\ell-1)}, \z_s^{(\ell)}, \z_t^{(\ell+1:L)}\Big] \in \Tilde{\mathcal{Z}}_\ell(\z_t^{(1:L)}).
\end{align*}
We additionally train an unconditional denoising network $p_\theta(\z_s^{(\ell)} \mid \z_t^{(1:L)})$ and then sample from the re-normalized distribution:
\begin{align}\label{eq:CBG}
    p^\gamma_{\phi, \theta}(\z_s^{(1:L)} \mid \z_t^{(1:L)}, y) = \prod_{\ell=1}^{L}\frac{p_\phi(y \mid \Tilde{\z}^{(1:L)})^\gamma p_\theta(\z_s^{(\ell)} \mid \z_t^{(1:L)})}{\sum_{\Tilde{\z}^{(1:L)}}p_\phi(y \mid \Tilde{\z}^{(1:L)})^\gamma p_\theta(\z_s^{(\ell)} \mid \z_t^{(1:L)})}.
\end{align}
Restricting the summation in the denominators of (\ref{eq:CBG}) to $\Tilde{\mathcal{Z}}_\ell(\z_t^{(1:L)})$ makes normalization tractable, as we are only summing over $\vocabsize$ terms.

We refer to our method as \textbf{D-CBG} for \textbf{D}iscrete \textbf{C}lassifier-\textbf{B}ased \textbf{G}uidance.
Our method can be thought of as an adaptation of the successful FUDGE \citep{yang2021fudge} approach, which guides AR generation, to discrete diffusion, similar to how NOS \citep{gruver2024protein} extended the AR guidance mechanism of PPLM \citep{dathathri2019plug} to diffusion models.

\textbf{First Order Approximation}~
While tractable, this formulation suffers from the drawback that at each denoising step we must perform $\gO(L\cdot \vocabsize)$ forward passes through the classifier model, which can quickly become impractical for larger vocabularies and longer sequences.
Similarly to \citet{grathwohl2021oops}, \citet{vignac2022digress}, 
% and \citet{nisonoff2024unlocking}, 
we can treat the classifier $p_\phi$ as a continuous function of the one-hot inputs $\tilde{\z}^{(1:L)} \in \mathbb{R}^{L \times \vocabsize}$ and use the first order Taylor approximation of $\log p_\phi$ to efficiently compute $p_{\phi}(y \mid \tilde{\z}^{(1:L)})$ with only a single forward and backward pass through the classifier model:
\begin{align}\label{eq:cbg_taylor}
    p_\phi(y \mid \tilde{\z}^{(1: L)}) &= \exp\left(\log\frac{p_\phi(y \mid \tilde{\z}^{(1:L)})}{p_\phi(y \mid \z_t^{(1:L)})} + \log p_\phi(y \mid \z_t^{(1:L)})\right) 
    \nonumber \\
    &\approx \exp\left((\tilde{\z}^{(1:L)} - \z_t^{(1:L)})^T\nabla_{\z_t^{(1:L)}}\log p_\phi(y \mid \z_t^{(1:L)}) + \log p_\phi(y \mid \z_t^{(1:L)})\right).
\end{align}

\section{Uniform Diffusion Language Models}\label{sec:method}
While masked diffusion models have demonstrated better language modeling compared to other discrete diffusion \citep{austin2021structured, lou2023discrete}, we argue that they are less amenable to guidance, since 
once a token is unmasked at some time $t$ it remains so for all $s < t$.
In contrast, for uniform noising, intermediate latents can be refined multiple times during the denoising process.
We therefore revisit categorical uniform noise discrete diffusion, where $\prior = \uniform := \vone / \vocabsize.$
Our aim is that by analyzing this class of diffusion models more carefully, we can reduce the gap to absorbing-state and yield performant models that are more easily steered by the guidance tools we developed above.

\subsection{Uniform Noise Diffusion}
% We begin by introducing the background and formulation for interpolating discrete diffusion models with uniform noise.
\textbf{Discrete Time Likelihood Bound}~
We start with the variational lower bound that is defined by a general diffusion process.
For discrete-time diffusion, i.e., some finite steps $T$, we define $t(i) = (i+1)/T$ and $s(i) = i / T$ for $i$ in $0, \ldots, T.$
Denoting the Kullback-Leibler divergence as $\KL$, the denoising network $p_\theta$ is trained to minimize a variational upper bound (NELBO), which is given by:
\begin{equation}\label{eq:nelbo}
    \E_q\Bigg[\underbrace{- \log p_\theta(
    \x | \z_{t(0)})}_{
    \begin{array}{c}\lossrecons\end{array}} + \underbrace{\sum_{i=1}^T \KL[q(\z_{s(i)} | \z_{t(i)}, \x) \| p_\theta(\z_{s(i)} | \z_{t(i)})]}_{
    \begin{array}{c}\lossdiff\end{array}}\Bigg]
    + \underbrace{\KL[q(\z_{t(T)} | \x) \| p_\theta(\z_{t(T)})]}_{
    \begin{array}{c}\lossprior\end{array}},
\end{equation}
When clear, we drop the explicit dependence of $t$ and $s$ on the discrete step $i$.

\textbf{Uniform Noise Forward Process}~
We focus on uniform noise diffusion using the interpolating discrete diffusion framework \citep{sahoo2024simple, zhao2024improving, zheng2023reparameterized}.
When letting $\prior = \uniform$, the input $\x$ transitions to a random state with some probability at each time step. 
Crucially, after $\x$ changes once, it can do so again.
When $\prior = \uniform$, the posterior from \eqref{eq:interpolating_posterior} becomes
\begin{equation}\label{eq:uniform_posterior}    
    q(\z_s \mid \z_t, \x)
    = \cat \left(
        \z_s;
        \frac{
            \vocabsize\at \z_t \odot \x + (\ats - \at)\z_t + (\as - \at)\x + \frac{(\as - \at)(1- \as)}{\vocabsize\as}\vone
            }
            {
             \vocabsize \at\langle \z_t, \x\rangle  + 1 - \at
            }
        \right)
\end{equation}

\textbf{Denoising Process}~
The optimal form for the reverse diffusion process $p_\theta$ matches (\ref{eq:uniform_posterior}): in fact setting $p_\theta$ to \eqref{eq:uniform_posterior} reduces the KL terms in \eqref{eq:nelbo} to zero. However, setting $p_\theta$ to exactly \eqref{eq:uniform_posterior} is not possible because it cannot be a function $\x$ (which $p_\theta$ is generating). 
Therefore, we introduce a predictive model $\x_\theta(\z_t, t): \onehots \times [0, 1] \rightarrow \Delta^N$ of the `clean' data given a noisy latent $\z_t$ at time $t$. We use $\x_\theta$ to parameterize the denoising process as $p_\theta(\z_s \mid \z_t) = q(\z_s \mid \z_t, \x = \x_\theta)$, yielding:
\begin{equation}\label{eq:uniform_posterior_pred}
    p_\theta(\z_s \mid \z_t)
    = \cat \left(
        \z_s;
        \frac{
            \vocabsize\at \z_t \odot \x_\theta + (\ats - \at)\z_t + (\as - \at)\x_\theta + \frac{(\as - \at)(1- \as)}{\vocabsize\as}\vone
            }
            {
             \vocabsize \at\langle \z_t, \x_\theta\rangle  + 1 - \at
            }
        \right),
\end{equation}
Note that this minimizes the $\lossdiff$ term in \eqref{eq:nelbo} precisely when $\x_\theta = \x,$ as desired.
To simplify notation, we omit below the explicit dependence of $\x_\theta$ on $t$.

\subsection{Improved Likelihood Bounds in Continuous Time}\label{subsec:udlm_elbo}
We now present our contribution to uniform noise discrete diffusion modeling.
We leverage the above formulation to develop an improved NELBO by taking $T\rightarrow\infty$ and analyzing each term $\lossrecons, \lossdiff, \lossprior$ in \eqref{eq:nelbo}.
This yields three improvements: (1) a simple and elegant closed-form expression for the NELBO that is easier to reason about; (2) an analytical reduction of $\lossrecons, \lossprior$ to zero, which tightens the NELBO; (3) a further tightening via the continuous-time extension of $\lossdiff$, as in \citet{kingma2021variational} and \citet{sahoo2024simple,sahoo2024mulan}.

\textbf{Prior Loss ($\lossprior$)}~
Given that we define our corruption process as an interpolation between clean data and a limiting distribution, for any noise schedules where $\alpha_{t(T)} = 0$, the distribution $q(\z_{t(T)}) = \prior$.
Therefore, we can simply define $p_\theta(\z_{t(T)}) = \prior$, and the KL divergence in $\lossprior$ evaluates to zero.

\textbf{Reconstruction Loss ($\lossrecons$)}~
For $\lossrecons$, if our noise schedule is such that  $T \rightarrow \infty \implies \alpha_{t(0)} = 1$ (i.e., $t(0)=1/T$), then the marginal $q(\z_{\frac{1}{T}} \mid \x) \rightarrow \cat(\z_{\frac{1}{T}}; \x)$.
That is, in the limit, the first latent vector is identically equal to the clean data.
We can thus parameterize our denoising network such that at time $t(0)$ the function simply copies its inputs: $\x_\theta(\z_{\frac{1}{T}}, 1/T) = \z_{\frac{1}{T}}.$
Additionally, we note that our choice of parameterization for $p_\theta$ implies that $p_\theta(\x \mid \z_{{\frac{1}{T}}}) = \x_\theta(\z_{{\frac{1}{T}}}, 1/T).$
Thus in the continuous time limit, we have:
\begin{equation*}
    \lim_{T\rightarrow\infty}\E_q[\lossrecons] = \lim_{T\rightarrow\infty}\E_q[\log p_\theta(\x\mid\z_{\frac{1}{T}})]
    = \lim_{T\rightarrow\infty}\E_q[\log(\langle\x, \x_\theta(\z_{\frac{1}{T}}, 1/T)\rangle)] = 0.
\end{equation*}

\textbf{Diffusion Loss ($\lossdiff$)}~
Turning finally to the diffusion loss term $\lossdiff$, we first define the shorthand $\KL[q_t||p_\theta] = \KL[[q(\z_{s} \mid \z_{t}, \x) || p_\theta(\z_{s}\mid \z_{t})]]$ and then re-write this loss term as an expectation over $t$ uniformly sampled from $1/T, 2/T, \ldots, T$:
\begin{equation}\label{eq:lossdiff_exp}
    \lossdiff = \sum_{i=1}^T\KL[q_{i/T}|| p_\theta]= T\cdot\E_{t\sim\{1/T, 2/T, \ldots, T\}}\KL[q_{t}|| p_\theta].
\end{equation}
Plugging in our expressions for the true and predicted posteriors from \eqref{eq:uniform_posterior} and \eqref{eq:uniform_posterior_pred} into 
\eqref{eq:lossdiff_exp}, then taking $T\rightarrow \infty$, we get (see \Appsecref{appsec:cont_uniform} for details):
\begin{equation}\label{eq:udlm_diff}
\begin{split}
    \lim_{T\rightarrow\infty}&\lossdiff =
    \lim_{T\rightarrow\infty}\E_{t\sim\{1/T, 2/T, \ldots, T\}}T\cdot\KL[q_t|| p_\theta] = \int_{t=0}^{t=1}\lim_{T\rightarrow\infty}T\cdot\KL[q_t||p_\theta]\mathrm{d}t
    \\
    &= \int_{t=0}^{t=1}\Bigg[
    \frac{\at'}{\vocabsize\at}\Bigg[\frac{\vocabsize}{\barx_i} - \frac{\vocabsize}{(\barx_\theta)_i} -\sum_{\substack{j \text{ s.t. }(\z_t)_j = 0}} \Bigg(\frac{\barx_j}{\barx_i}\Bigg) \log\Bigg[\Bigg(\frac{(\barx_\theta)_i\cdot\barx_j}{(\barx_\theta)_j\cdot\barx_i}\Bigg)\Bigg]
    \Bigg]\Bigg]\mathrm{d}t,
\end{split}
\end{equation}
where $\x_j$ denotes the $j^{\text{th}}$ index of a vector $\x$, $\barx = \vocabsize\at\x + (1 - \at)\vone$, $\barx_\theta = \vocabsize\at\x_\theta + (1 - \at)\vone$, and we define $i = \arg\max_{j\in[\vocabsize]} (\z_t)_j$ to be the non-zero entry of $\z_t$.

Combining our arguments regarding $\lossprior$ and $\lossrecons$ with \eqref{eq:udlm_diff}, yields our final tight bound:
\begin{equation}\label{eq:udlm_elbo_token}
    \losscont = \int_{t=0}^{t=1}\E_q\Bigg[
    \frac{\at'}{\vocabsize\at}\Bigg[\frac{\vocabsize}{\barx_i} - \frac{\vocabsize}{(\barx_\theta)_i} -\sum_{\substack{j \text{ s.t. }(\z_t)_j = 0}} \Bigg(\frac{\barx_j}{\barx_i}\Bigg) \log\Bigg[\Bigg(\frac{(\barx_\theta)_i\cdot\barx_j}{(\barx_\theta)_j\cdot\barx_i}\Bigg)\Bigg]
    \Bigg]\Bigg]\mathrm{d}t.
\end{equation}

\textbf{Extension to Sequences}~
Extending training with \eqref{eq:udlm_elbo_token} from $\x \in \onehots$ to sequences $\x^{(1:L)}\in\onehots^L,$ we make the assumption that the denoising process factorizes independently across tokens when conditioned on a sequence of noisy latents $\z_t^{(1:L)}.$
In this case, we use a single model $\x^{(\ell)}_\theta(\z_t^{(1:L)}, t)$ for predicting each token $\ell \in \{1, \ldots, L\}$ in a sequence, and we train with the sequence-level objective:
\begin{equation}
    \losscont
    = \int_{t=0}^{t=1}\E_q \sum_\ell\Bigg[
    \frac{\at'}{\vocabsize\at}\Bigg[\frac{\vocabsize}{\barx_i^{(\ell)}} - \frac{\vocabsize}{(\barx_\theta^{(\ell)})_i} -\sum_{\substack{j \text{ s.t. }(\z_t^{(\ell)})_j = 0}} \Bigg(\frac{\barx_j^{(\ell)}}{\barx_i^{(\ell)}}\Bigg) \log\Bigg[\Bigg(\frac{(\barx_\theta^{(\ell)})_i\cdot\barx_j^{(\ell)}}{(\barx_\theta^{(\ell)})_j\cdot\barx_i^{(\ell)}}\Bigg)\Bigg]
    \Bigg]\Bigg]\mathrm{d}t.
\end{equation}
We dub models trained with our refined objective \textbf{U}niform \textbf{D}iffusion \textbf{L}anguage \textbf{M}odels (UDLM).

\section{Experiments}\label{sec:exp}
% \vspace{-1em}
\textbf{Datasets}~
For our language modeling experiments, we examine several discrete domains:
reference genomes from tens species (Species10),
the QM9 small molecule dataset \citep{ruddigkeit2012enumeration, ramakrishnan2014quantum}, where molecules are represented by SMILES strings \citep{weininger1988smiles},  CIFAR10 discretized images \citep{krizhevsky2009learning},
and three NLP datasets consisting of text8 \citep{text8}, Amazon Review \citep{mcauley2013hidden, zhang2015character}, and the one billion words dataset (LM1B; \citet{chelba2014billion}).
These datasets cover a range of domains and vocabularies of varying sizes (see \Tabref{tab:ppl}).
For guidance experiments, we explore species-specific sequence generation, molecular property maximization, and class-conditional image generation.

\subsection{Language Modeling with Uniform Noise Discrete Diffusion}\label{subsec:exp_lm}
Our language modeling experiments show that (1) contrary to a widely-held belief, \textbf{uniform noise diffusion can attain state-of-the-art performance} on small vocabulary datasets (\Tabref{tab:ppl}), and that (2) our \textbf{\method~are state-of-the-art among uniform noise diffusion} (Tables \ref{tab:text8} and \ref{tab:lm1b}).

\begin{wraptable}[15]{r}{7.2cm}
\vspace{-1em}
\small
\caption{\method~performs best with smaller vocabs.
Perplexity ($\downarrow$) on various datasets.
Best values are \textbf{bolded.}
$^{\text{*}}$ indicates values reported from early stopping on the validation set; otherwise validation performance at the end of training is used.
$^\dagger$From \cite{sahoo2024simple}.
$^\$$From \citet{lou2023discrete}.
}
\label{tab:ppl}
\vspace{-1em}
\begin{center}
\begin{tabular}{lcccc}
\toprule
 & $|$Vocab.$|$ & AR & MDLM & \method \\%& \%$\Delta$ MDLM - UDLM \\
\midrule
% \multicolumn{5}{l}{\textit{Scientific applications}} \\
Species10 & 12 & \bf{2.88} & 3.17$_{\leq}$ & 3.15$_{\leq}$ \\
QM9$^{\text{*}}$ & 40 & 2.19 & 2.12$_{\leq}$ &\bf{2.02}$_{\leq}$ \\
% \multicolumn{5}{l}{\textit{Images}} \\
CIFAR10 & 256 & - & \bf{9.14}$_{\leq}$ & 11.21$_{\leq}$ \\
% \multicolumn{5}{l}{\textit{NLP}} \\
text8 & 35 & \bf{2.35}$^{\$}$ & 2.62$_{\leq}$ & 2.71$_{\leq}$ 
\\
Amazon$^{\text{*}}$ & 30,522 & \bf{21.67} & 24.93$_{\leq}$ & 27.27$_{\leq}$ \\
LM1B & 30,522 & \bf{22.32}$^{\dagger}$ & 27.04$^{\dagger}_{\leq}$ & 31.28$_{\leq}$ \\
\bottomrule
\end{tabular}
\end{center}
\end{wraptable}

In \Tabref{tab:ppl}, despite previous evidence indicating that absorbing-state discrete diffusion greatly outperforms uniform noise, we find a more nuanced story.
Namely, for smaller vocabulary regimes, the gap between MDLM and \method~is negligible, with \method~even outperforming absorbing-state on certain datasets.
Even within a single domain, we find a trend between vocabulary size and performance gap, with the text8 results of \method~on par with MDLM, and a more persistent gap for larger vocabularies.
Intuitively, for larger vocabulary regimes, uniform noise diffusion models need to predict over a combinatorially larger set of potential clean data sequences compared to absorbing-state; smaller vocabularies reduce this complexity.

\begin{wrapfigure}[12]{l}{0.38\textwidth}
\vspace{-2em}
    \begin{center}
        \includegraphics[
        width=\linewidth, 
        ]{images/amazon_ppl.pdf}
    \end{center}
\vspace{-1.5em}
\caption{$T \rightarrow \infty$ improves
validation PPL ($\downarrow$) on Amazon dataset.}
\label{fig:amazon}
\end{wrapfigure}
Secondly, despite the persisting gap between absorbing and uniform noise discrete diffusion for larger vocabulary NLP datasets, we note that models trained with \method~close this gap, as evidenced in Tables \ref{tab:text8} and \ref{tab:lm1b}, where we show that \method~attains the best reported uniform noise discrete diffusion language modeling performance.

\textbf{Ablating Continuous Time Formulation}~
In \Figref{fig:amazon}, we see the effect of increasing $T$ and using our continuous time NELBO.
The curves represent validation perplexity on the Amazon Polarity dataset at various training steps, and we observe that increasing $T$ indeed improves language modeling with \method, i.e., $T=\infty$, performing best.

\begin{figure}
\begin{minipage}[t]{0.49\linewidth}
\centering
\small
\captionof{table}{\method~outperforms other uniform discrete diffusion on text8.
Best value is \textbf{bolded} \& best uniform diffusion value is \underline{underlined}.
$^\dagger$From \cite{lou2023discrete}.
$^\text{*}$From \cite{shi2024simplified}.
}
\label{tab:text8}
\begin{tabular}{llc}
    \toprule
    & Method & BPC ($\downarrow$) \\
    \midrule
    \multicolumn{3}{l}{\textit{Autoregressive}} \\
    & IAF/SCF$^\dagger$ 
    \scriptsize{\citep{ziegler2019latent}}
    % [\citenum{ziegler2019latent}] 
    & 1.88 \\
    & Argmax Flow$^\dagger$ 
    \scriptsize{\citep{hoogeboom2021argmax}}
    % [\citenum{hoogeboom2021argmax}]
    & 1.39 \\
    & Discrete Flow$^\dagger$ 
    \scriptsize{\citep{tran2019discrete}}
    % [\citenum{tran2019discrete}]
    & \textbf{1.23} \\
    & Autoregressive$^\dagger$ & \textbf{1.23} \\
    \multicolumn{3}{l}{\textit{Non-autoregressive}} \\
    & Mult. Diffusion$^\dagger$ 
    % [\citenum{hoogeboom2021argmax}]
    \scriptsize{\citep{hoogeboom2021argmax}}
    & 1.72$_\leq$ \\
    & MAC$^\dagger$
    \scriptsize{\citep{shih2022training}}
    % [\citenum{shih2022training}]
    & 1.40$_\leq$ \\
    & BFN$^\dagger$ 
    \scriptsize{\citep{graves2023bayesian}}
    % [\citenum{graves2023bayesian}]
    & 1.41$_\leq$ \\
    & D3PM Absorb$^\dagger$ 
    \scriptsize{\citep{austin2021structured}}
    % [\citenum{austin2021structured}] 
    & 1.45$_\leq$ \\
    & SEDD Absorb$^\dagger$ 
    \scriptsize{\citep{lou2023discrete}}
    % [\citenum{lou2023discrete}]
    & 1.39$_\leq$ \\
    & MDLM
    \scriptsize{\citep{sahoo2024simple}}
    % [\citenum{sahoo2024simple}]
    \textit{(Retrained)}
    & 1.38$_\leq$ \\
    & MD4$^\text{*}$
    \scriptsize{\citep{shi2024simplified}}
    % [\citenum{shi2024simplified}]
    & 1.37$_\leq$ \\
    & GenMD4$^*$
    \scriptsize{\citep{shi2024simplified}}
    % [\citenum{shi2024simplified}]
    & 1.34$_\leq$ \\
    \midrule
    \multicolumn{3}{l}{\textit{Discrete Uniform Diffusion}} \\
    & D3PM Uniform$^\dagger$ 
    \scriptsize{\citep{austin2021structured}}
    % [\citenum{austin2021structured}] 
    & 1.61$_\leq$ \\
    & SEDD Uniform$^\dagger$ 
    \scriptsize{\citep{lou2023discrete}}
    % [\citenum{lou2023discrete}]
    & 1.47$_\leq$ \\
    & UDLM (\textit{Ours}) & \underline{1.44}$_\leq$\\
    \bottomrule
\end{tabular}
\end{minipage}
\hfill
\begin{minipage}[t]{0.46\linewidth}
\centering
\small
\captionof{table}{\method~outperforms other uniform discrete diffusion on LM1B.
Best value is \textbf{bolded} \& best discrete uniform diffusion value is \underline{underlined}.
$^\dagger$From \cite{sahoo2024simple}.
$^\text{*}$From \cite{lou2023discrete}.
$^\$$From \cite{austin2021structured}.
}
\label{tab:lm1b}
\begin{tabular}{llc}
    \toprule
    & Method & PPL ($\downarrow$) \\
    \midrule
    \multicolumn{3}{l}{\textit{Autoregressive}$^\dagger$} \\
    & Transformer-X
    \scriptsize{\citep{dai2019transformer}}
    % [\citenum{dai2019transformer}]
    & 23.5 \\
    & OmniNetT 
    \scriptsize{\citep{tay2021omninet}}
    % [\citenum{tay2021omninet}]
    & \bf{21.5} \\
    & Transformer & 22.32 \\
    \multicolumn{3}{l}{\textit{Diffusion}$^\dagger$} \\
    & BERT-Mouth 
    \scriptsize{\citep{wang2019bert}}
    % [\citenum{wang2019bert}]
    & 142.89$_\leq$ \\
    & D3PM Absorb 
    \scriptsize{\citep{austin2021structured}}
    % [\citenum{austin2021structured}] 
    & 77.50$_\leq$ \\
    & Diffusion-LM 
    \scriptsize{\citep{li2022diffusion}}
    % [\citenum{li2022diffusion}]
    & 118.62$_\leq$ \\
    & DiffusionBert 
    \scriptsize{\citep{wang2019bert}}
    % [\citenum{wang2019bert}] 
    & 63.78$_\leq$ \\
    & SEDD Absorb 
    \scriptsize{\citep{lou2023discrete}}
    % [\citenum{lou2023discrete}]
    & 32.79$_\leq$ \\
    & MDLM
    \scriptsize{\citep{sahoo2024simple}}
    % [\citenum{sahoo2024simple}]
    & 27.04$_\leq$ \\
    \midrule
    \multicolumn{3}{l}{\textit{Discrete Uniform Diffusion}} \\
    & D3PM Uniform$^\text{\$}$ 
    \scriptsize{\citep{austin2021structured}}
    % [\citenum{austin2021structured}] 
    & 137.9$_\leq$ \\
    & SEDD Uniform$^\text{*}$ 
    \scriptsize{\citep{lou2023discrete}}
    % [\citenum{lou2023discrete}]
    & 40.25$_\leq$ \\
    & UDLM (\textit{Ours}) & \underline{31.28}$_\leq$ \\
    \bottomrule
\end{tabular}
\end{minipage}
\end{figure}

\subsection{Guided Discrete Diffusion}\label{subsec:results_guidance}
Our guidance results indicate that (1) \textbf{classifier-free guidance is more useful when paired with diffusion models compared to AR} (\Tabref{tab:dna_species}) and that (2) \textbf{our proposed \cbg~is the best classifier-based method} for discrete guidance, especially when combined with \method~(\Tabref{tab:qm9_qed} \& \Figref{fig:qm9_guidance}).

\textbf{Baselines}~
For guidance experiments, our primary baseline is the dominant AR approach.
% Our goal is to show that with similar, even slightly worse, language modeling performance, discrete diffusion models equipped with classifier-based and classifier-free guidance can outperform AR models on controlled generation.
We compare to three flavors of guided AR.
The first is applying \cfg~to AR models.
We also use the established control mechanisms of Plug-and-play language models (PPLM; \citet{dathathri2019plug}) and FUDGE \citep{yang2021fudge}.
To demonstrate the better performance of our \cbg~method, we compare to \citet{gruver2024protein} (NOS), an extension of PPLM to discrete diffusion.

\textbf{Hyperparameters}~
For both \cfg~and \cbg, we vary the strength of the $\gamma$ parameter.
Although the original FUDGE formulation simply uses $\gamma=1$, we also perform a search for this baseline.
For PPLM and NOS, we vary the parameters of the Langevin sampling that updates models' hidden representation, i.e., number of update steps $n$, step size $\eta$, and fluency KL weight $\gamma_{kl}$ (see \cite{dathathri2019plug} and \cite{gruver2024protein} for more details).
For all methods, we display the best performing hyperparameter configuration in the main table results, and defer the the full sweep to \Appsecref{appsec:guidance_ablation}.

\begin{wraptable}{r}{7cm}
\vspace{-1em}
\small
\caption{Diffusion decoding with \cfg~is more controllable than AR for genomic sequences.
Mean $\pm$ standard deviation reported from five random seeds.
Best values are \textbf{bolded}.
}
\label{tab:dna_species}
\vspace{-1em}
\begin{center}
\begin{tabular}{lcccc}
\toprule
Model & \makecell{\cfg\\$\gamma$} & \makecell{Disc.\\AUROC ($\downarrow$)} & F1 ($\uparrow$) \\
\midrule
Random & \textendash & 1.00 & 0.07 \\
AR & 1 & 0.53$_{\pm \text{0.075}}$ & 0.87$_{\pm\text{0.011}}$ \\
AR & 2 & 0.90$_{\pm\text{0.051}}$ & 0.81$_{\pm\text{0.009}}$ \\
AR & 3 & 0.97$_{\pm\text{0.018}}$ & 0.74$_{\pm\text{0.022}}$ \\
MDLM & 1 & \bf0.51$_{\pm\text{0.059}}$ & 0.88$_{\pm\text{0.011}}$ \\
MDLM & 2 & 0.74$_{\pm\text{0.037}}$ & 0.91$_{\pm\text{0.009}}$ \\
MDLM & 3 & 0.93$_{\pm\text{0.031}}$ & 0.78$_{\pm\text{0.007}}$ \\
UDLM & 1 & 0.52$_{\pm\text{0.044}}$ & 0.91$_{\pm\text{0.01}}$ \\
UDLM & 2 & 0.61$_{\pm\text{0.043}}$ & 0.93$_{\pm\text{0.009}}$ \\
UDLM & 3 & 0.87$_{\pm\text{0.051}}$ & \bf0.94$_{\pm\text{0.007}}$ \\
\bottomrule
\end{tabular}
\end{center}
\vspace{-1em}
\end{wraptable}
\textbf{Species-specific Genome Generation}~
For genomic sequences, we evaluate \cfg~with diffusion compared to with AR.
We train on sequences of 32,768 nucleotides, using base-pair level tokenization and conditionally generate 64 sequences for each class.
% As quality measures for each species class, we compute the Jensen-Shannon (JS) distance between the $k$-mer frequencies of the generated sequences and those from the validation set, and report the mean JS across species (weighted by species frequency in the dataset), with smaller values indicating better $k$-mer distributional overlap between the ground truth and generated sequences \citep{sarkar2024designing}.
% Additionally, 
We train a small classifier to distinguish between generated and validation set sequences and report the area under the receiver operator curve for this classifier (Disc. AUROC).
Values closer to 0.5 indicate that the classifier is unable to distinguish between synthetic and true sequences \citep{sarkar2024designing}.
To measure the controllability, we train a separate classifier on the Species10 dataset and measure macro F1 score of this `oracle' classifier on the generated sequences.
Results of this experiment are presented in \Tabref{tab:dna_species}.
For reference, we also provide metrics for randomly generating sequences with nucleotide frequencies proportional to species representation in the data.
We find that both MDLM and \method~are able to better generate sequences that match the desired control parameter, with higher F1 scores relative to AR.
Moreover, \method~is able to outperform MDLM in satisfying this control.
Importantly, we find that only \method~is amenable to increasing the guidance parameter $\gamma$, where its metrics improves while AR and MDLM metrics degrade.
Finally, of note, the diffusion model generation for this experiment is accomplished with far fewer function evaluations compared to AR.
Whereas AR must decode each of the 32,768 tokens, because MDLM and \method~can decode multiple tokens in parallel, we generate with $T = 512.$

\begin{table}[t!]
% \vspace{-2em}
\small
\caption{Guidance with discrete diffusion models better balances the generation of valid and novel molecules with maximizing the property of interest, drug likeness (QED) / ring count, compared to AR.
Validity, novelty, and mean QED / ring count for novel sequences are measured for generated sequences from each method.
Mean $\pm$ standard deviation reported from five random seeds.
Best values are \textbf{bolded}.
}
\vspace{-1em}
\label{tab:qm9_qed}
\begin{center}
\begin{tabular}{lllcccccc}
\toprule
& & & \multicolumn{3}{c}{\textit{QED}} & \multicolumn{3}{c}{\textit{Ring Count}}
\\
& Method & Guidance & Valid ($\uparrow$) & Novel ($\uparrow$) & Mean ($\uparrow$) & Valid ($\uparrow$) & Novel ($\uparrow$) & Mean ($\uparrow$) \\
\midrule
\multicolumn{2}{l}{Original Data} & \makecell[c]{\textendash} & 133k & 133k & 0.47 & 133k & 133k & 1.74\\
\midrule
\multicolumn{6}{l}{\textit{Classifier-free}}\\
& AR & \cfg & 946.4$_{\pm\text{9.0}}$ & 79.4$_{\pm\text{6.4}}$ & 0.60$_{\pm\text{0.00}}$ & 441.4$_{\pm\text{11.1}}$ & 77.8$_{\pm\text{2.3}}$ & 4.83$_{\pm\text{0.08}}$ \\
& MDLM & \cfg & 317.4$_{\pm\text{11.5}}$ & \bf{95.8}$_{\pm\text{9.0}}$  & 0.60$_{\pm\text{0.01}}$  & 90.0$_{\pm\text{8.2}}$ 	& 60.0$_{\pm\text{7.5}}$  & \bf{5.26}$_{\pm\text{0.15}}$ \\
& \method & \cfg & \bf{1013.6}$_{\pm\text{2.5}}$ & 64.0$_{\pm\text{5.1}}$ & \bf{0.62}$_{\pm\text{0.00}}$  & 998.2$_{\pm\text{4.5}}$  & 216.2$_{\pm\text{13.0}}$  & 4.88$_{\pm\text{0.04}}$ \\
\midrule
\multicolumn{6}{l}{\textit{Classifier-based}}\\
& AR & FUDGE & 924.4$_{\pm\text{12.1}}$ & 53.0$_{\pm\text{3.5}}$ & 0.61$_{\pm\text{0.00}}$  & 281.2$_{\pm\text{5.0}}$  & 103.6$_{\pm\text{4.8}}$ & 4.89$_{\pm\text{0.10}}$ \\
& AR & PPLM & 1007.2$_{\pm\text{2.3}}$ & 142.0$_{\pm\text{14.2}}$ &	0.45$_{\pm\text{0.00}}$ & \bf{1009.8}$_{\pm\text{1.9}}$ & 140.6$_{\pm\text{16.5}}$ & 1.92$_{\pm\text{0.09}}$ \\
& MDLM & \cbg & 417.6$_{\pm\text{19.7}}$ & 116.6$_{\pm\text{8.9}}$ & 0.58$_{\pm\text{0.00}}$ & 113.0$_{\pm\text{8.8}}$ & 85.6$_{\pm\text{8.8}}$ & 4.75$_{\pm\text{0.23}}$ \\
& MDLM & NOS & 506.0$_{\pm\text{29.5}}$ & \bf{240.4}$_{\pm\text{16.0}}$ & 0.45$_{\pm\text{0.01}}$  & 247.8$_{\pm\text{14.0}}$ & 193.6$_{\pm\text{13.1}}$ & 3.51$_{\pm\text{0.22}}$ \\
& \method & \cbg & 994.8$_{\pm\text{2.9}}$ & 63.8$_{\pm\text{8.1}}$ & 0.61$_{\pm\text{0.00}}$ & 897.2$_{\pm\text{11.6}}$ & \bf{432.0}$_{\pm\text{19.1}}$ & 4.84$_{\pm\text{0.02}}$ \\
& \method & NOS & 547.8$_{\pm\text{10.6}}$ & 158.8$_{\pm\text{11.0}}$ & 0.47$_{\pm\text{0.00}}$ & 573.8$_{\pm\text{13.0}}$ & 244.4$_{\pm\text{13.1}}$ & 3.96$_{\pm\text{0.07}}$ \\
\bottomrule
\end{tabular}
\end{center}
\end{table}

\begin{wraptable}[14]{l}{6cm}
    \small
    \centering
    \caption{Guidance improves FID \& IS on CIFAR10.
    Finite- (D3PM) vs. continuous-time (MDLM / UDLM).
    % Guidance using \cfg~($\gamma=4$).
    Best values are \textbf{bolded}.
    $^\dagger$From \cite{austin2021structured}.}
    \label{tab:cifar10}    
    \begin{tabular}{lccc}
    \toprule
     & FID ($\downarrow$) & IS ($\uparrow$)\\
     \midrule
     D3PM Absorb$^\dagger$ & 41.28 & 6.26 
     & \\
    MDLM & 33.75 & 6.74 \\ 
    MDLM \cfg & \bf15.56 & \bf9.02\\    
    \midrule
    D3PM Uniform$^\dagger$ & 51.27 & 5.99 \\
    \method & 33.65 & 6.86 \\
    \method~\cfg & 23.21 & 8.66 \\
    \bottomrule
    \end{tabular}
\end{wraptable}
\textbf{Molecular Property Maximization}~
For QM9, we investigate novel generation of sequences that maximize either drug-likeness (QED; \citet{bickerton2012quantifying}) or number of rings present in the molecule.
We only report values for which we generated at least 50 novel sequences (out of 1,024).
In \Tabref{tab:qm9_qed}, for \cfg, we see that all methods perform comparably when maximizing drug likeness, but that MDLM and \method~are better suited for the structural property of ring count.
For classifier-based guidance mechanisms, we again find that both QED and ring count, \cbg~with discrete diffusion best trades-off maximizing these properties while generating valid and novel sequences compared to AR or diffusion with NOS.

% \begin{table}[ht!]
% \small
% \caption{Guidance with discrete diffusion models better balances the generation of valid and novel molecules with maximizing the property of interest, drug likeness (QED), compared to AR.
% Validity, novelty, and mean QED for novel sequences are measured for generated sequences from each method.
% Mean $\pm$ standard deviation reported from repeated generation of 1,024 sequences using five different random seeds.
% Best mean values are \textbf{bolded}.
% ``no approx.'' subscript for \cbg~implies first order approximation was \textit{not} used.
% }
% \label{tab:qm9_qed}
% \begin{center}
% \begin{tabular}{lllccc}
% \toprule
% & Method & Guidance & Valid ($\uparrow$) & Novel ($\uparrow$) & QED / Mean ($\uparrow$) \\
% \midrule
% \multicolumn{2}{l}{Original Data} & \makecell[c]{-} & 133k & 133k & 0.47\\
% \midrule
% \multicolumn{6}{l}{\textit{Classifier-free}}\\
% & AR & \cfg$_{\gamma=3}$ & 946.4 \scriptsize{$\pm$ 8.99} & 79.4 \scriptsize{$\pm$ 6.35} & 0.60 \scriptsize{$\pm$ 0.00} \\
% & MDLM & \cfg$_{\gamma=3}$ & 317.4 \scriptsize{$\pm$ 11.5} & \bf{95.8} \scriptsize{$\pm$ 9.04} & 0.60 \scriptsize{$\pm$ 0.01} \\
% & \method & \cfg$_{\gamma=5}$ & \bf{1013.6} \scriptsize{$\pm$ 2.51} & 64.0 \scriptsize{$\pm$ 5.1} & \bf{0.62} \scriptsize{$\pm$ 0.00} \\
% \midrule
% \multicolumn{6}{l}{\textit{Classifier-based}}\\
% & AR & FUDGE$_{\gamma=7}$ & 924.4 \scriptsize{$\pm$ 12.05} & 53.0 \scriptsize{$\pm$ 3.54}	& 0.61 \scriptsize{$\pm$ 0.00} \\
% & AR & PPLM$_{n=30, \eta=0.1}$ & 1007.2 \scriptsize{$\pm$ 2.28} & 142.0 \scriptsize{$\pm$ 14.2} &	0.45 \scriptsize{$\pm$ 0.00} \\
% & MDLM & \cbg$_{\gamma=3, \text{no approx.}}$ & 417.6 \scriptsize{$\pm$ 19.69} & 116.6 \scriptsize{$\pm$ 8.91} & 0.58 \scriptsize{$\pm$ 0.00} \\
% & MDLM & NOS$_{n=5, \eta=0.1, \gamma_{kl}=0.01}$ & 506.0 \scriptsize{$\pm$ 29.45} & \bf{240.4} \scriptsize{$\pm$ 16.01} & 0.45 \scriptsize{$\pm$ 0.01} \\
% & \method & \cbg$_{\gamma=10, \text{no approx.}}$ & 994.8 \scriptsize{$\pm$ 2.86} & 63.8 \scriptsize{$\pm$ 8.11} & 0.61 \scriptsize{$\pm$ 0.00} \\
% & \method & NOS$_{n=10, \eta=5, \gamma_{kl}=0.01}$ & 547.8 \scriptsize{$\pm$ 10.57} & 158.8 \scriptsize{$\pm$ 10.99} & 0.47 \scriptsize{$\pm$ 0.00} \\
% \bottomrule
% \end{tabular}
% \end{center}
% \end{table}

\begin{figure}
    % \vspace{-1em}
    \centering
    \includegraphics[width=0.48\linewidth]{images/qed_cbg.pdf}
    \includegraphics[width=0.48\linewidth]{images/ring_count_cfg.pdf}
    \caption{
    Diffusion models extend the steer-ability Pareto frontier.
    \textit{(Left)} 
    D-CBG outperforms FUDGE classifier guidance when maximizing drug-likeness (QED).
    % We display $\gamma$ values for which 50 or more of the generated sequences are novel.
    % Only UDLM can accommodate larger $\gamma$ values that enable better molecule property maximization.
    \textit{(Right)}
    % Comparing D-CFG for diffusion and AR models for maximizing ring-count in the QM9 dataset.
    \cfg~with diffusion better trades-off novel generation and ring-count maximization compared to AR.
    }
    \label{fig:qm9_guidance}
    % \vspace{-1em}
\end{figure}

% \begin{table}[ht!]
% \small
% \caption{Guidance with UDLM best balances the generation of valid and novel molecules with maximizing the property of interest, ring count, compared to AR and MDLM.
% Validity, novelty, and mean ring count for novel sequences are measured for generated sequences from each method.
% Mean $\pm$ standard deviation reported from repeated generation of 1,024 sequences using five different random seeds.
% Best mean values are \textbf{bolded}.
% ``1$^\text{st}$ order'' subscript for \cbg~implies first order approximation was used and
% ``no approx.'' subscript implies approximation was \textit{not} used.
% }
% \label{tab:qm9_ring}
% \begin{center}
% \begin{tabular}{lllccc}
% \toprule
% & Method & Guidance & Valid ($\uparrow$) & Novel ($\uparrow$) & Ring Count Mean ($\uparrow$) \\
% \midrule
% \multicolumn{2}{l}{Original Data} & \makecell[c]{-} & 133k & 133k & 1.74\\
% \midrule
% \multicolumn{6}{l}{\textit{Classifier-free}} \\
% & AR & \cfg$_{\gamma=5}$ & 441.4 \scriptsize{$\pm$ 11.08} & 77.8 \scriptsize{$\pm$ 2.28} & 4.83 \scriptsize{$\pm$ 0.08} \\
% & MDLM & \cfg$_{\gamma=5}$ & 90.0 \scriptsize{$\pm$ 8.15}	& 60.0 \scriptsize{$\pm$ 7.52} & \bf{5.26} \scriptsize{$\pm$ 0.15} \\
% & \method & \cfg$_{\gamma=3}$ & 998.2 \scriptsize{$\pm$ 4.49} & 216.2 \scriptsize{$\pm$ 12.99} & 4.88 \scriptsize{$\pm$ 0.04} \\
% \midrule
% \multicolumn{6}{l}{\textit{Classifier-based}} \\
% & AR & FUDGE$_{\gamma=5}$ & 281.2 \scriptsize{$\pm$ 4.97} & 103.6 \scriptsize{$\pm$ 4.83}	& 4.89 \scriptsize{$\pm$ 0.1} \\
% & AR & PPLM$_{n=10, \eta=0.1}$ & \bf{1009.8} \scriptsize{$\pm$ 1.92} & 140.6 \scriptsize{$\pm$ 16.47} & 1.92 \scriptsize{$\pm$ 0.09} \\
% & MDLM & \cbg$_{\gamma=10, \text{1}^{\text{st}}\text{ order}}$ (Ours) & 113.0 \scriptsize{$\pm$ 8.75} & 85.6 \scriptsize{$\pm$ 8.82} & 4.75 \scriptsize{$\pm$ 0.23} \\
% & MDLM & NOS$_{n=5, \eta=5, \gamma_{kl}=0.01}$ & 247.8 \scriptsize{$\pm$ 13.97} & 193.6 \scriptsize{$\pm$ 13.13} & 3.51 \scriptsize{$\pm$ 0.22} \\
% & \method & \cbg$_{\gamma=8, \text{no approx.}}$ (Ours) & 897.2 \scriptsize{$\pm$ 11.58} & \bf{432.0} \scriptsize{$\pm$ 19.07}	& 4.84 \scriptsize{$\pm$ 0.02} \\
% & \method & NOS$_{n=5, \eta=5, \gamma_{kl}=0.01}$ & 573.8 \scriptsize{$\pm$ 13.03} & 244.4 \scriptsize{$\pm$ 13.05} &	3.96 \scriptsize{$\pm$ 0.07}
%  \\
% \bottomrule
% \end{tabular}
% \end{center}
% \end{table}

\textbf{Class-conditional Image Generation}~
In \Tabref{tab:cifar10}, we see that both MDLM and \method~outperform finite-time counterparts (in the form of D3PM \citep{austin2021structured}) with improved image quality metrics of Fr\'echet inception distance (FID; \citet{heusel2017gans}) and Inception Score (IS; \citet{salimans2016improved}).
This is especially true when we add guidance using \cfg.
MDLM / \method~samples are generated with $T=1000$.
% In \Figref{fig:cifar10}, we show several conditionally generated images.

\textbf{Ablation: Faster Sampling}~
In \Tabref{tab:cifar10_T_ablation}, we explore using faster inference settings (smaller $T$) where diffusion models predict multiple pixels in parallel.
These results focus on the CIFAR10 dataset; see \Appsecref{appsubsec:T_ablation} for additional experiments on other datasets.
MDLM's performance deteriorates with smaller $T$, whereas \method~is robust to this setting.
This validates a key motivation behind \method: in settings where MDLM `locks in' certain predictions that it cannot change, \method~is more resilient given that all tokens can change throughout the decoding process.

% \begin{figure}[ht]
% \vspace{-10em}
% \centering
% \begin{minipage}[c]{0.55\linewidth}
%     \small
%     \centering
%     \captionof{table}{Guidance improves quality on discretized CIFAR10.
%     FID and IS for finite- (D3PM) and continuous-time (MDLM / UDLM) discrete diffusion models. Guidance using \cfg~($\gamma=4$).
%     Best values are \textbf{bolded}.
%     $^\dagger$From \cite{austin2021structured}.}
%     \label{tab:cifar10}
%     % \begin{adjustbox}{width=\linewidth}
%     \begin{tabular}{lccc}
%     \toprule
%      & FID ($\downarrow$) & IS ($\uparrow$)\\
%      \midrule
%      D3PM Absorb$^\dagger$ & 41.28 & 6.26 
%      & \\
%     MDLM & 33.75 & 6.74 \\ 
%     MDLM \cfg & \bf15.56 & \bf9.02\\    
%     \midrule
%     D3PM Uniform$^\dagger$ & 51.27 & 5.99 \\
%     \method & 33.65 & 6.86 \\
%     \method~\cfg & 23.21 & 8.66 \\
%     \bottomrule
%     \end{tabular}
%     % \end{adjustbox}
% \end{minipage}
% \hfill
% \begin{minipage}[c]{0.41\linewidth}
%     % % \vspace{-1em}
%     % \includegraphics[width=\linewidth]{images/cifar10.pdf}
%     % \vspace{-1em}
%     % \caption{CIFAR10 conditional generation.}
%     % \label{fig:cifar10}
%     \small
%     \centering
%     \captionof{table}{UDLM is robust to faster sampling.
%     Images sampled from a conditional model (\cfg$_{\gamma=1}$) trained on CIFAR10.
%     F1 metric from a separate classifier trained to identify class label used for generation.
%     Best metric per $T$ is \textbf{bolded}.
%     }
%     \label{tab:cifar10_T_ablation}
%     \begin{tabular}{llccc}
%         \toprule
%          & Model & FID ($\downarrow$) & IS ($\uparrow$) & F1 ($\uparrow$) \\
%          \midrule
%          \multicolumn{5}{l}{$T=128$} \\
%          & MDLM & 64.09 & 5.81 & 0.63\\
%          & UDLM & \bf{30.48} & \bf{7.30} & \bf{0.80}\\
%          % \midrule
%          % \multicolumn{5}{l}{$T=256$} \\
%          % & MDLM & 41.53 & 6.51 & 0.70\\
%          % & UDLM & \bf{28.30} & \bf{7.40} & \bf{0.81}\\
%          % % \midrule
%          % \multicolumn{5}{l}{$T=512$} \\
%          % & MDLM & 31.97 & 6.89 & \bf{0.74}\\
%          % & UDLM & \bf{27.12} & \bf{7.43} & \bf{0.74}\\
%          % \midrule
%          \multicolumn{5}{l}{$T=1024$} \\
%          & MDLM & 27.94 & 7.14 & \bf{0.81} \\
%          & UDLM & \bf{26.70} & \bf{7.43} & \bf{0.81} \\
%          \bottomrule
%     \end{tabular}
% \end{minipage}
% \end{figure}

\section{Related Works, Discussion, and Conclusion}\label{sec:related}
\begin{wraptable}{r}{5.5cm}
    \vspace{-1.1em}
    \small
    \centering
    \caption{UDLM is robust to faster sampling.
    CIFAR10 images sampled from a conditional model (\cfg$_{\gamma=1}$).
    % F1 metric from a separate classifier trained to identify class label used for generation.
    Best values per $T$ are \textbf{bolded}.
    }
    \label{tab:cifar10_T_ablation}
    \begin{tabular}{llcc}
        \toprule
         & Model & FID ($\downarrow$) & IS ($\uparrow$) \\
         \midrule
         \multicolumn{4}{l}{$T=128$} \\
         & MDLM & 64.09 & 5.81\\
         & UDLM & \bf{30.48} & \bf{7.30}\\
         % \midrule
         % \multicolumn{5}{l}{$T=256$} \\
         % & MDLM & 41.53 & 6.51 & 0.70\\
         % & UDLM & \bf{28.30} & \bf{7.40} & \bf{0.81}\\
         % % \midrule
         % \multicolumn{5}{l}{$T=512$} \\
         % & MDLM & 31.97 & 6.89 & \bf{0.74}\\
         % & UDLM & \bf{27.12} & \bf{7.43} & \bf{0.74}\\
         % \midrule
         \multicolumn{4}{l}{$T=1024$} \\
         & MDLM & 27.94 & 7.14\\
         & UDLM & \bf{26.70} & \bf{7.43} \\
         \bottomrule
    \end{tabular}
\end{wraptable}
\textbf{Discrete Diffusion}~
Recent works have examined interpolating discrete diffusion \citep{ou2024your, sahoo2024simple, shi2024simplified, zhao2024improving, zheng2023reparameterized}, a special case of the general framework from D3PM \citep{austin2021structured}.
Our \method~method most closely aligns to the state-of-the-art discrete diffusion of \cite{ou2024your}, \cite{sahoo2024simple}, and \cite{shi2024simplified} that focus on absorbing state diffusion.
Similar to these works, we provide a continuous time NELBO leading to performance gains, but we focus on uniform noise.
Of note, other extensions of D3PM that do not start from the variational perspective,  instead rely on the formalisms of continuous time Markov chains (CTMC) \citep{campbell2022continuous} and concrete score matching \citep{lou2023discrete}, but are less performant than works such as MDLM and our method.
Also stemming from CTMC are extensions of flow-based \citep{lipman2022flow} approaches to discrete data \citep{campbell2024generative, gat2024discrete}.
In \Appsecref{appsec:ctmc}, we discuss connections of our work to CTMC in more detail.

\textbf{Guidance}~
Leveraging continuous embeddings of discrete data, Diffusion-LM \citep{li2022diffusion} uses Langevin sampling with CBG.
Similarly, SSD-LM \citep{han2022ssd} perform Gaussian noising on the logits of a bi-directional model, which they combine with pre-trained classifiers to perform CBG with Langevin dynamics.
LD4LG \citep{lovelace2024latent} implement CFG on continuous embeddings.
\cite{stark2024dirichlet} use flow-matching on the simplex to perform CFG and CBG on continuous representations.
\citet{wang2023infodiffusion} perform guidance using auxiliary semantic latent variables.
In contrast, to these continuous formulations for discrete data, our work adapts guidance mechanisms directly to the discrete domain.

% \begin{wraptable}{r}{6cm}
%     \vspace{-1.5em}
%     \centering
%     \caption{UDLM is robust to faster sampling.
%     Images sampled from a conditional model (\cfg$_{\gamma=1}$) trained on CIFAR10.
%     F1 metric from a separate classifier trained to identify class label used for generation.
%     Best metric per $T$ is \textbf{bolded}.
%     }
%     \begin{tabular}{llccc}
%         \toprule
%          & Model & FID ($\downarrow$) & IS ($\uparrow$) & F1 ($\uparrow$) \\
%          \midrule
%          \multicolumn{5}{l}{$T=128$} \\
%          & MDLM & 64.09 & 5.81 & 0.63\\
%          & UDLM & \bf{30.48} & \bf{7.30} & \bf{0.80}\\
%          % \midrule
%          % \multicolumn{5}{l}{$T=256$} \\
%          % & MDLM & 41.53 & 6.51 & 0.70\\
%          % & UDLM & \bf{28.30} & \bf{7.40} & \bf{0.81}\\
%          % % \midrule
%          % \multicolumn{5}{l}{$T=512$} \\
%          % & MDLM & 31.97 & 6.89 & \bf{0.74}\\
%          % & UDLM & \bf{27.12} & \bf{7.43} & \bf{0.74}\\
%          % \midrule
%          \multicolumn{5}{l}{$T=1024$} \\
%          & MDLM & 27.94 & 7.14 & \bf{0.81} \\
%          & UDLM & \bf{26.70} & \bf{7.43} & \bf{0.81} \\
%          \bottomrule
%     \end{tabular}
%     \label{tab:cifar10_T_ablation}
% \end{wraptable}

FUDGE \citep{yang2021fudge} can be viewed as the AR analog to our \cbg.
DiGress \citep{vignac2022digress} uses a similar first-order approximation to resolve the intractability of the normalizing constant.
In our work, we derive a tractable expression for classifier-based guidance and simply use the Taylor approximation to speed up computation in large sequence length and vocabulary size regimes.
% \citet{li2024derivative} take a reinforcement-learning lens to this problem and use a value function that estimates the expected future reward given the current noisy discrete latent and use this estimation to weight the denoising distribution of a pre-trained unconditional discrete diffusion model.
\citet{sanchez2023stay} derives an equivalent formulation for \cfg~and use it to better enforce AR models' adherence to prefix prompts.
FreeGress \citep{ninniri2024classifier} also offer a comparable method to \cfg, but focus on graph diffusion models.
Most similar to our method, is the concurrent work of \citet{nisonoff2024unlocking}, which derives classifier-based and classifier-free guidance for discrete diffusion and flow models.
However, this work is highly tailored to models that leverage the formalism of CTMC, and guidance is applied to the rate matrices.
See \Appsecref{appsec:ctmc_guidance} for details on relating our guidance methods to \citet{nisonoff2024unlocking}.

\textbf{Conclusion}~
In search of a more controllable diffusion process, in this work, we derived a tight variational bound for uniform noise discrete diffusion, closing the gap to state-of-the-art absorbing-state diffusion models.
We also highlighted that contrary to previous findings, in small vocabulary regimes, uniform noise is on par or better than absorbing state.
We then demonstrated that straightforward adaptations of classifier-based and classifier-free guidance can offer improved guided generation relative to AR models.
We found that with classifier-free mechanisms, diffusion models are more amenable to control without sacrificing quality of generated sequences.
We also demonstrated that our classifier-based method is better than previous ones for both AR and diffusion models.

\section*{Acknowledgments and Disclosure of Funding}
We thank Nate Gruver for useful discussions regarding NOS. 
This work was supported by NSF CAREER grants (\#2145577 and \#2037519) and an NIH MIRA grant (\#1R35GM151243-01).


\bibliography{refs}
\bibliographystyle{iclr2025_conference}

\newpage
\appendix
\tableofcontents
\input{appendix}

\end{document}
